\section{Life in Kholinar}
\label{sec:kholinarLife}

Start the adventure by reading the following:

\quoteBox{
	It is a few days after Lightday - the Weeping still continues to pour endless rains over the city of Kholinar. In the puddles of water gathered on the streets Rinspren stretch upwards towards the clouds like a blue candles. \\ \\
	Few people can be found on the streets in this weather - most often those are parshmen workers hauling baskets and boxes, never complaining on the weather conditions.
}

This adventure doesn't have any particular scenes or events to start the story. In the first scenes, focus on introducing the characters and bringing them together - your goal should be to get ready to move towards \fullref{sec:sweetHome}. This section will provide several examples of scenes you can use. In each scene you play in this chapter, you should remember to show:

\begin{itemize}
	\item The NPCs you want to introduce
	\item Parshmen working in the background
	\item Food being scarce and expensive
	\item Populace growing unhappy with the queen's rule
\end{itemize}

\reasonBox{
	People being unsatisfied with the current rule and lack of food will trigger major events in \nameref{chap:chapter2}. Use this chapter to paint the picture for your players.
	\\ \\
	Parshmen will play major role in \nameref{chap:chapter3}, so start foreshadowing them right away. There will be less opportunity for that in \nameref{chap:chapter2}, so don't wait with that.
}

\newpage
\subsubsection{\npcTroublemaker{} the Troublemaker}

All example scenes for this part of the adventure feature \npcTroublemaker{}. It is assumed she knows at least some people in the party (as a distant family or just a friend).

\roleplayBox{\npcTroublemaker{}}{
	\roleplayTips{Impulsive, reckless, enthusiastic}
	{Help \npcGrandma{} survive in starving city}
	{Young and thin, with light eyes an strands of hair unable to keep in order}
	{\npcTroublemaker{} is the embodiment of the proverb "Curiosity killed the skyeel". Trouble seem to follow her, often as a direct consequence of her impulsive actions. This doesn't seem to bother her, as she can always find something fun in every scenario.}
}

\reasonBox{
	The role of \npcTroublemaker{} in the story is to ensure the players will see the most important events of the plot. Being reckless and impulsive, she will rush towards any danger without thinking. If the PCs wish to protect her, they will need to step out of their comfort zone and risk some unsafe actions.
}


\subsection{Messing with the Wrong People}

Start this scene by reading the following:

\quoteBox{
	You stroll on the streets of Kholinar, when suddenly you hear a racket behind you. You turn to see a thin teenage girl running away from a few thugs chasing her. \\ \\
	"Look, I'm sorry I stole your food, but I'm not giving it back!", she shouts to one of them, then she quickly adds looking at you: "A little help, perhaps?"
}

The thugs (use \enemyRioter{} stat block) are chasing \npcTroublemaker{}. There are as many of them as the PCs (at least two). If the characters intervene, the thugs will fight until there is only one of them left, at which point he will escape. 

\npcTroublemaker{} will not participate in the fight, but she will thank the characters for their help.

\reasonBox{
	If possible, weave \npcTroublemaker{} into your characters' backstories. She will appear quite a lot in the first chapters, often driving the plot forward.
}



\subsection{Hungry Axehounds}

When you start this scene, read the following:

\quoteBox{
	You see a teenage girl surrounded by a pack of axehounds, clicking aggressively at her. The girl holds some package above her head. \\ \\
	"It's my food, guys!", she pleads with the pack. "Go buy your own!"
}

\npcTroublemaker{} is being surrounded by \enemyAxehound{}s (two Axehounds per PC). If attacked, the pack will fight to the end.

The characters may also attempt an \skillIntimiation{} test or a \skillLeadership{} test against respective Defenses of an \enemyAxehound{} to scare them or command them away. Any test is made with disadvantage unless the PC uses some food before to befriend the animals. 




\subsection{Soaking Wet}

When you start this scene, read the following:

\quoteBox{
	You see a group of parshmen doing laundry in a clean waters of the Weeping. Nearby, a young woman observes their efforts. \\ \\
	"Do you think they have a havah in there?", \npcTroublemaker{} asks. "I would need one, my old one is... well... well, I need a new one!"
}

\npcTroublemaker{} attempts to find and steal a havah dress from the parshmen. The PCs can help her in this attempt, or they may try to persuade her to do better.

\subsubsection{Helping \npcTroublemaker{}}

If the characters wish to help \npcTroublemaker{}, you can run the rest of the scene as an endeavor.
\pursuitReq{5}{3}

The party must use \skillPerception{} to find the dress, \skillThievery{} and \skillStealth{} to take it without being noticed and perhaps some \skillIntimiation{} or \skillDeception{} if they are caught by the workers.

\subsubsection{Stopping \npcTroublemaker{}}

If the characters wish to stop \npcTroublemaker{}, you can run the rest of the scene as a conversation.
For the purpose of the scene, you can assume \npcTroublemaker{} is an \enemyExpert{}.
\conversationDefence{\npcTroublemaker{}}{13}{13}{4}
She leads a reckless life full of impulsive choices, but she is not a bad person.

\conversationEnd{\npcTroublemaker{}}{1}



\subsection{Chull-sliding}

Start this scene with the following:

\quoteBox{
	You enter a marketplace, closed for the trade at this hour. There, a chull caravan is resting after a long travel. \\ \\
	"Hey, look, a real chulls!", \npcTroublemaker{} says. "Have you ever played a chull-sliding? I always wanted to try!", she adds before running to the caravan.
}

\npcTroublemaker{} will climb the biggest chull to slide down on its shell. The PCs can join her. Unless they finish the game soon enough, read the following:

\quoteBox{
	Suddenly a merchant appears in the nearest building's window. "Thieves! They steal my chulls! Call the guards!"
}

A merchant - \enemyExpert{} - enters the scene. The PCs can quickly enter a conversation to calm him down. 
\conversationDefence{The merchant}{13}{13}{4}
He is very protective of the caravan and finds this kind of games foolish.

\conversationEnd{the merchant}{1}

If the party fails or ignores attempts to calm the merchant down, four \enemyWallGuard{}s will soon appear. The party can fight them or flee the scene. 

If they choose the latter, you can run it as a pursuit endeavor with a Starting Distance of 2 and an Escape Distance of 4. If the guards catch the party, they will attempt to arrest them, which probably will start a combat encounter.



\subsection{Grateful Invitation}

If the party helps \npcTroublemaker{} to deal with her troubles, she will be really grateful for the aid. As a thank you, she will invite the PCs to her home, so they could eat something and warm themselves in a dry house. Proceed to \fullref{sec:sweetHome}.